% Two-column preview approximating published JCIM layout
% Compile: pdflatex preview_twocol && bibtex preview_twocol && pdflatex preview_twocol && pdflatex preview_twocol
%
% NOT for submission — use benchmark_paper_jcim.tex (achemso) for that.
% This previews the published two-column format.
\documentclass[10pt,twocolumn,letterpaper]{article}

% ===== Page geometry (ACS two-column) =====
\usepackage[top=0.75in, bottom=0.75in, left=0.63in, right=0.63in, columnsep=0.25in]{geometry}

% ===== Packages (matching JCIM version) =====
\usepackage{amsmath,amssymb,amsfonts}
\usepackage{graphicx}
\usepackage{booktabs}
\usepackage{subcaption}
\usepackage{siunitx}
\usepackage{adjustbox}
\usepackage{soul}
\usepackage{tikz}
\usetikzlibrary{arrows.meta, positioning, shapes, calc, fit, decorations.pathreplacing}
\usepackage{longtable}
\usepackage{multirow}
\usepackage[dvipsnames]{xcolor}
\usepackage{placeins}
\usepackage[numbers,super,sort&compress]{natbib}
\usepackage[colorlinks,citecolor=blue!60!black,linkcolor=blue!60!black,urlcolor=blue!60!black]{hyperref}
\usepackage{mathptmx}            % Times Roman (ACS house style)
\usepackage[scaled=0.92]{helvet} % Helvetica sans-serif
\usepackage{courier}
\usepackage[font=small,labelfont=bf]{caption}
\usepackage{titlesec}
\usepackage{float}

% ===== Section formatting =====
\titleformat{\section}{\normalfont\bfseries\large}{\thesection.}{0.5em}{}
\titleformat{\subsection}{\normalfont\bfseries}{\thesubsection.}{0.5em}{}
\titlespacing*{\section}{0pt}{1.2ex plus 0.4ex}{0.8ex plus 0.2ex}
\titlespacing*{\subsection}{0pt}{1.0ex plus 0.3ex}{0.6ex plus 0.1ex}

% ===== Custom Commands =====
\newcommand{\vect}[1]{\mathbf{#1}}
\newcommand{\matr}[1]{\mathbf{#1}}

% ===== Compact float spacing =====
\setlength{\textfloatsep}{8pt plus 2pt minus 2pt}
\setlength{\floatsep}{6pt plus 2pt minus 2pt}
\setlength{\intextsep}{6pt plus 2pt minus 2pt}
\setlength{\abovecaptionskip}{4pt}
\setlength{\belowcaptionskip}{2pt}
\setlength{\dbltextfloatsep}{8pt plus 2pt minus 2pt}
\setlength{\dblfloatsep}{6pt plus 2pt minus 2pt}

% ===== Title =====
\title{\Large\bfseries When Do Simple Models Win? Machine Learning Architectures for UV Absorption Prediction and Insights for General Molecular Property Prediction}

\author{%
Umesh Arampath,$^{a,b,*}$ Bryant Pero,$^{b}$ David Stewart,$^{a}$ and David Demirjian$^{b}$\\[6pt]
{\small\itshape $^{a}$Department of Mathematics, University of Iowa, Iowa City, IA, USA}\\
{\small\itshape $^{b}$Midwest Bioprocessing Center, IL, USA}\\[3pt]
{\small $^{*}$E-mail: uarampath@mwbioprocessing.com}
}

\date{}

\begin{document}

\twocolumn[
  \maketitle
  \begin{@twocolumnfalse}
  \begin{abstract}
\noindent
What determines which machine learning architecture wins for molecular property prediction?
We address this question through a controlled comparison of five models spanning four architecture families: fingerprint-based ensembles (Random Forest, XGBoost), a directed message-passing graph neural network (GNN), a bidirectional gated recurrent network (BiGRU), and a pretrained Transformer, using UV absorption wavelength ($\lambda_{\max}$) prediction as a testbed across 18,755 solute--solvent pairs with stratified 5-fold cross-validation and two external datasets totaling over 40,000 molecules.
The answer depends on the task.
For screening within known chemical space, RF with Morgan fingerprints is optimal: it matches the GNN (RMSE = $31.34 \pm 1.82$ vs.\ $31.69 \pm 3.18$~nm, $p = 0.77$), trains in 15 minutes on a CPU with no hyperparameter tuning, and outperforms all sequence-based models.
For exploring novel scaffolds, deep learning is superior: experimental validation on 16 novel UV-absorbing compounds shows that the GNN (MAE = 26.2~nm), Transformer (26.3~nm), and an optimized BiGRU (28.6~nm; 3-layer, 256-unit architecture identified via Bayesian hyperparameter search) all outperform RF (38.5~nm), as learned representations generalize where fixed fingerprints cannot.
In either case, encoding solvent identity improves all models: a simple SMILES/fingerprint concatenation strategy yields 6--8\% RMSE reductions across fingerprint, graph, and recurrent architectures ($p < 0.02$), without domain-specific descriptor engineering.
Interpretability analysis confirms that both RF feature importance and BiGRU gradient saliency independently highlight the same chromophore motifs, and that architectures with local inductive bias systematically outperform global-attention Transformers on this locally determined property.
These findings provide general guidance for molecular property prediction: the largest or most complex model is not always the best choice. Instead, optimal model selection depends on the locality of the target property and the relationship between target compounds and available training data. When the target property is governed by local molecular features and the prediction task is interpolation within known chemical space, simple fingerprint-based models offer an efficient and accurate alternative to deep learning.
  \end{abstract}
  \vspace{0.5em}
  \noindent\textbf{Keywords:} UV absorption $\cdot$ molecular property prediction $\cdot$ deep learning $\cdot$ random forest $\cdot$ Morgan fingerprints $\cdot$ graph neural network $\cdot$ directed message passing $\cdot$ pretrained Transformer $\cdot$ solvent effects
  \vspace{1em}
  \end{@twocolumnfalse}
]

% ===== Full paper body (auto-extracted from benchmark_paper_jcim.tex) =====
% Abstract is in the wrapper above; body starts at Introduction
\input{preview_body}

% ===== Bibliography =====
\bibliographystyle{achemso}
\bibliography{Proposal}

\end{document}
